\section[R transport and missingness contracts]{R transport and missingness: preserve the population, the missingness law, and the analysis target}
\label{sec:exec25-rtm}

A bank moving a successful servicing experiment to a different customer population faces two related but distinct questions. Which customers should the experimental effect describe, and what can be learned when outcomes or adjustment variables are unavailable? A transport estimator changes the population over which an identified response is averaged. An imputation model supplies a distribution for unobserved values under declared restrictions. Neither operation automatically supplies the assumptions needed by the other. This distinction governs the proposed R interface and its comparison with generalize, transportr, causaltransportr, mice, smcfcs, and causalOT. The inspected releases are pinned before the client cutoff, with precise identities and primary implementation anchors. The comparisons concern those selected implementations, not every release. No downloaded package was executed.\citep{exec25_rtm_generalize,exec25_rtm_transportr,exec25_rtm_causaltransportr,exec25_rtm_mice,exec25_rtm_smcfcs,exec25_rtm_causalOT}

\subsection{The population contract precedes the estimator}

Let $S=1$ identify the source trial and $S=0$ the target sample, $A$ be a binary treatment, $X$ the pre-treatment covariates, and $Y$ the outcome in declared units. Write $\rho(x)=P(S=1\mid X=x)$ and $e(x)=P(A=1\mid X=x,S=1)$. Under consistency, source treatment exchangeability, overlap, and equality of the relevant conditional potential-outcome means across populations, the target contrast is
\[
 \tau_0=E\{m_1(X)-m_0(X)\mid S=0\},\qquad
 m_a(x)=E(Y\mid A=a,X=x,S=1).
\]
These are restrictions on the causal and sampling problem, not consequences of fitting a selection classifier. A corresponding augmented population functional is
\begin{align*}
 \tau_0={1\over P(S=0)}E\bigg[&(1-S)\{m_1(X)-m_0(X)\}\\
 &+S{1-\rho(X)\over\rho(X)}
 \left\{{A\over e(X)}(Y-m_1(X))
       -{1-A\over1-e(X)}(Y-m_0(X))\right\}\bigg].
\end{align*}
With the true selection and treatment probabilities, conditioning on $X,S=1$ reduces the residual term to the difference between the true and working response contrasts, multiplied by the target covariate density. It therefore repairs a misspecified response model in this population calculation. With the correct response model, that residual has conditional expectation zero. These algebraic statements do not establish finite-sample coverage after nuisance fitting, cluster sampling, or clipping. The authored finite-law reference checks the identity separately from those extensions.\citep{exec25_rtm_choice}

The generalize weighting delegate distinguishes disjoint source and target samples, using $(1-\rho)/\rho$, from a nested-population mode using $1/\rho$. It removes rows missing selection covariates, normalizes target survey weights by their target mean, and optionally caps source weights at an empirical percentile. Its random-forest path expands survey rows using rounded-up counts and fixes a sampling-size argument at 454. An adapter must record the retained population, sampling interpretation, and delegate arguments rather than treat all settings as equivalent weighting.\citep{exec25_rtm_generalize}

The package also illustrates why an output named TATE is insufficient to identify its scale. The inspected weighting routine uses a Gaussian model when the outcome has more than two observed values, but uses a quasibinomial fit and exponentiates its treatment coefficient when there are exactly two. The latter quantity is an odds ratio in the applicable binary-logit model, not a risk difference. For risks $0.1$ and $0.2$, the difference is $0.1$, while the odds ratio is $9/4$. Multiplying the latter by customer counts and a per-event value would be dimensionally wrong. The result envelope must carry the estimand, coding, link, contrast, and scale separately, and reject ambiguous two-valued numeric outcomes.\citep{exec25_rtm_generalize,exec25_rtm_choice}

Population trimming is equally consequential. generalize forces its trimming option in the selected TMLE dispatch. Its helper retains factor levels and numerical coordinate ranges represented in the trial. This does not establish joint support: source points $(0,0)$ and $(1,1)$ admit target point $(0,1)$ under both marginal ranges even though the joint point has no source support. Conversely, excluding customers changes the target to its retained subset. Both populations and the excluded mass must be reported; a smaller interval after restriction is not improved evidence for every original customer.\citep{exec25_rtm_generalize,exec25_rtm_choice}

\subsection{Read the score and its complete delegate chain}

transportr separates a task, nuisance fits, cross-fitting, and an influence-function object. Its standard score uses source inverse odds and an average response difference over target observations. The task exposes population roles, observation indicators, and cluster identifiers. Nevertheless, its default one-fold setting uses the same observations for training and validation. More folds activate separate fold construction, with a cluster-aware branch when identifiers are supplied. Cluster-aware splitting alone does not derive cluster-aware uncertainty. The selected influence routine hands row contributions to the external ife package without a cluster identifier in that call; the complete uncertainty delegate remains a qualification requirement. The documented sampling-weight argument is currently unused.\citep{exec25_rtm_transportr}

The standard score also needs a bounded missing-outcome contract. Its outcome regression is trained on observed source rows, whereas its residual helper replaces every missing product with zero. The displayed standard influence function contains no explicit inverse observation-probability factor. This requires a separately justified missingness argument; an observation-indicator argument does not establish general censoring robustness. Zero times an unavailable target outcome is a structural omission, while a missing source outcome or failed nuisance prediction is a different event. The adapter must distinguish them instead of translating all missing products into successful zeros.\citep{exec25_rtm_transportr}

A separate iterated-outcome delegate gives another reason to qualify an exact branch. Both its treatment-one and treatment-zero fits select the variable named \texttt{.outcome\_eps1}, although the task also carries \texttt{.outcome\_eps0}. This source-level routing concern requires a public fixture with distinct arm-specific pseudo-outcomes. It is not a measured failure rate and does not invalidate an unrelated standard-ATE branch. Reuse must be specified by task, estimator and delegate.\citep{exec25_rtm_transportr}

causaltransportr supplies explicit generalize, transport and in-sample targets with augmented weighting, unaugmented weighting, outcome modeling and calibration alternatives. Its nuisance code fits $S$ and constructs source-membership probabilities. In the inspected no-surrogate transport branch, the augmented score uses $S\rho/(1-\rho)$, while the unaugmented branch uses $S(1-\rho)/\rho$. These are different weights under that definition of $\rho$. To expose the consequence without a fitted package, take $P(S=1)=4/5$, independent randomized treatment with probability $1/2$, $Y=A$, and a working outcome regression identically zero. The target effect is one. The expected augmented residual with the displayed reversed odds, divided by $P(S=0)$, is 16; correct inverse odds give one. This is an exact population witness for the selected formula under supplied nuisance values, not a claim that every estimator or native run produces 16.\citep{exec25_rtm_causaltransportr,exec25_rtm_choice}

The optional normalization needs its own contract. Selected branches divide the residual matrix by the column sums of treatment probabilities, which are not sums of the displayed inverse probability weights. With ten rows and probability $1/2$ throughout, the divisor is five. The native parity fixture must check matrix recycling, dimensions and the intended normalized score. The summary delegate divides sums by a possibly substituted target denominator while calculating standard deviations over the supplied contribution rows. Its uncertainty formula must be derived for the same normalization and target-sample design. Source inspection identifies the questions; it does not supply the variance theorem.\citep{exec25_rtm_causaltransportr,exec25_rtm_choice}

\subsection{Imputation preserves a model, not missing causal information}

mice offers chained equations, blocks, predictor matrices, passive transformations and multiple-imputation analysis. Its selected sampler forms fitting rows from complete cases, observed outcomes, and the complement of the ignore indicator. Its documentation states that some multivariate delegates do not honor ignore and can require a complete predictor space incompatible with arbitrary where masks. A training/test boundary must therefore be verified for the chosen imputer. A global flag cannot certify that validation outcomes have been withheld from every delegate. Its monotone one-iteration shortcut also edits the predictor matrix without checking that the data actually have the assumed pattern.\citep{exec25_rtm_mice}

Predictive mean matching returns observed donor values. That can preserve support and avoid implausible draws, but it cannot generate an unobserved value outside donor support: with donors zero and one, a missing value of two cannot be produced by that rule. This is a limitation of the mechanism, not a claim that extrapolation is always desirable. The product should expose donor count, exclusions, matching rule and missingness sensitivity, rather than describe every filled table as recovered truth.\citep{exec25_rtm_mice,exec25_rtm_choice}

The selected MNAR normal imputer makes this boundary explicit by adding a user-specified unidentifiable-model term to a normal draw. Those sensitivity coefficients are supplied restrictions, not quantities learned from unobserved outcomes. In a simple mean calculation with missing fraction $0.3$, changing the missing-group mean by an elicited $\delta=2$ changes the overall mean by $0.6$. The output should be a sensitivity curve indexed by declared coefficients and their interpretation. Its spread is not automatically a confidence interval for an identified effect. A chosen coefficient cannot certify the bank's reporting or assignment mechanism.\citep{exec25_rtm_mice,exec25_rtm_choice}

This distinction survives pooling. mice explicitly warns against pooling completed datasets before fitting the scientific model. It combines parameter estimates and variances, preferring a robust standard-error field when supplied by the model's tidy method. For estimates $(0,2,4)$ with within-imputation variances all one, the between-imputation variance is four. Its missing-data rule gives total variance $1+(1+1/3)4=19/3$, whereas its partially synthetic-data rule gives $1+4/3=7/3$. Neither validates an incompatible imputation model, a misspecified causal target, or an unsupported complete-data variance. The interface must preserve pooling rule, degrees of freedom, parameter identity, and covariance requirements. Scholarly references in package comments are not newly cited here without genuine local originals.\citep{exec25_rtm_mice,exec25_rtm_choice}

smcfcs addresses compatibility with a declared substantive model. Its discrete covariate update multiplies covariate-model probabilities by the substantive outcome likelihood and normalizes the products. With prior $P(X=1)=1/2$, $P(Y=1\mid X=1)=0.8$, and $P(Y=1\mid X=0)=0.2$, an observed $Y=1$ changes the compatible imputation probability to $0.8$. Drawing from the marginal probability $1/2$ would discard that information. Passive transformations are updated within the procedure, so interactions and nonlinear terms must remain consistent through the calculations. Compatibility with a chosen model still does not establish its scientific truth.\citep{exec25_rtm_smcfcs,exec25_rtm_choice}

For continuous proposals the inspected implementation first makes up to 24 vectorized rejection attempts, then uses a finite per-subject batch. If every candidate is rejected, it increments a failure counter; the previously stored value remains, and the routine later warns while returning imputed datasets. Under an illustrative constant acceptance probability $0.1$ and batch limit two, rejecting all 26 proposals has probability $0.9^{26}$. This is not an observed package rate or a universal rejection law. It shows why suppressing warnings erases a scientific failure condition. The wrapper must retain warnings, affected records, variables, iterations and limits, and refuse clean completion until a qualified policy handles unresolved draws. Increasing the limit is a possible computational remedy, not a proof of compatibility or convergence.\citep{exec25_rtm_smcfcs,exec25_rtm_choice}

\subsection{Optimal transport and adoption decisions}

causalOT uses optimal-transport methods to construct causal weights; this mathematical transport does not automatically transport experimental effects to a new population. The selected API distinguishes ATT, ATC and ATE, carries initial sample weights, and forms weighted contrasts with optional outcome-model augmentation. Its numeric effective-sample-size method evaluates $1/\sum_i w_i^2$ under a documented sum-to-one convention. Raw weights $(1,3)$ therefore give $0.1$, while normalized weights $(1/4,3/4)$ give $1.6$. Effective sample size must carry its normalization convention and cannot certify joint support or valid effect uncertainty.\citep{exec25_rtm_causalOT,exec25_rtm_choice}

Reuse decisions are conditional and specific. generalize is a candidate for population weighting and diagnostics after target, scale and trimming qualification. transportr offers task and nuisance boundaries but needs branch-specific observation, routing and uncertainty qualification. Direct admission of the selected causaltransportr score is deferred pending correction or a justified interpretation of its odds and normalization, followed by public parity tests. mice is a candidate imputation and pooling engine under a per-method contract; smcfcs is a candidate where its substantive-model and rejection-failure conditions match the request. causalOT is a candidate weighting component, not a substitute for a transportability argument. A port must reproduce scientific decisions and failure states as well as successful arithmetic; direct adoption must retain versioned runtime behavior.\citep{exec25_rtm_choice}

Source declarations distinguish MIT for generalize and causaltransportr, GPL version three or later for transportr, GPL version two or later for mice, and the selected GPL version-three declarations for smcfcs and causalOT. These do not settle complete dependency, native-library, service or redistribution rights. Pinned bundles retain identity, hashes, repository blob checks and inspected anchors. Four bounded discovery queries preserve returned candidates and omissions; no exhaustive ecosystem or saturation claim follows. Papers mentioned by packages remain discovery leads until genuine originals and technical anchors are available.\citep{exec25_rtm_generalize,exec25_rtm_transportr,exec25_rtm_causaltransportr,exec25_rtm_mice,exec25_rtm_smcfcs,exec25_rtm_causalOT,exec25_rtm_choice}

Public qualification can precede expensive bank experiments. Fixtures should freeze selection law, intervention, population weights, outcome scale, nuisance values, joint support, imputation mask and sampling unit; compare exact scores before fitting learners; then test fitted nuisance functions, dependence, failure propagation and costs separately. CLI and MCP requests should serialize these fields and return typed refusals for unsupported combinations. In the hypothetical odds witness, a discrepancy of 15 outcome units across 10,000 customers valued at USD10 per unit changes the arithmetic by USD1.5 million. This motivates prototype tests, not a measured loss or ROI claim. The immediate route is concrete while authentic bank truth and adoption remain later gates: qualify the population and missingness contract before interpreting a completed dataset as a completed causal argument.\citep{exec25_rtm_choice}
